\begin{recipe}{Roze koeken}
\kicker[Bakken,Zoet,Kinderen]{Bakken · zoet · kinderen}
\index[register]{Roze koeken}
\index[register]{Frambozen}
\index[register]{Citroen}

\meta{45 minuten + afkoeltijd \dvd Ca.\ 12 stuks \dvd Oven 175\,°C}

\lettrine{R}{oze} koeken zijn de Nederlandse versie van een petit four. Het
muffinblik dient als mal, en een bakplaat erop tijdens het bakken zorgt voor
een platte bovenkant om in te dippen.

\blockrule{Ingrediënten}
\begin{ingredients}
  \ing{100 g}{boter, op kamertemperatuur (ca.\ 91\%)}
  \ing{100 g}{witte basterdsuiker (ca.\ 91\%)}
  \ing{2}{eieren}
  \ing{1 tl}{vanille-extract}
  \ing{110 g}{patentbloem}
  \ing{1}{citroen, gewassen, voor de rasp}\index[register]{Citroen}
  \ing{1/4 tl}{zout}
\end{ingredients}

\blockrule{Glazuur}
\begin{ingredients}
  \ing{100 g}{frambozen (diepvries)}\index[register]{Frambozen}
  \ing{200 g}{poedersuiker}
\end{ingredients}

\blockrule{Bereiding}
\tip{Laat de koekjes na het bakken eerst even in de vorm rusten voor je ze eruit haalt
en ondersteboven op een rekje legt. Zo blijft de onderkant mooi vlak voor het glazuur.}
\begin{steps}
  \item Verwarm de oven voor op 175\,°C (boven- en onderwarmte). Vet een muffinblik in met boter en bestrooi
        licht met bloem.
  \item Klop de boter en basterdsuiker luchtig met een garde.
  \item Voeg de eieren één voor één toe en meng goed na elk ei.
  \item Voeg de citroenrasp en het vanille-extract toe en klop nog een minuut door.
  \item Zeef het zout en de bloem door het mengsel en spatel voorzichtig door.
  \item Doe het beslag in een spuitzak en vul het muffinblik.
  \item Leg een bakplaat op het muffinblik zodat de koekjes niet kunnen rijzen. Bak
        10 minuten met de bakplaat erop, daarna nog 6 minuten zonder.
  \item Laat de koekjes 5 minuten rusten, haal ze uit de vorm en leg ze
        ondersteboven op een rekje tot ze volledig zijn afgekoeld.
  \item Ontdooi de frambozen en wrijf ze door een bolzeef zodat de pitjes achterblijven.
  \item Meng de helft van de frambozenpuree met de poedersuiker tot een glad
        glazuur en voeg meer puree toe als het te dik is.
  \item Doop de platte kant van elke koek in het glazuur en laat opdrogen op een rekje.
\end{steps}

\heroimagefade[recto]{images/roze-koeken}
\end{recipe}
